% =============================================================================
% Kapitel 3: Hauptteil / Analyse
% -----------------------------------------------------------------------------
% Der Kern der Arbeit. Hier wird die Forschungsfrage bearbeitet. Pro
% Unterkapitel eine zentrale Aussage, im Dreischritt: Begriff/Gegenstand klären
% -- entfalten -- Folgerung ziehen und überleiten.
% =============================================================================
\section{Analyse}
\label{sec:hauptteil}

Aufbauend auf den Grundlagen aus Kapitel~\ref{sec:grundlagen} bearbeitet dieses
Kapitel die Forschungsfrage. Es entwickelt das Argument schrittweise und führt
die Teilergebnisse zusammen.

% --- 3.1 ---
\subsection{Erster Analyseschritt}
\label{subsec:analyse1}

Jeder Absatz trägt genau einen Gedanken und ist belegt. Mehrere Quellen lassen
sich kombinieren.\autocite[Vgl.][S.~22]{mustermann2024}\autocite[Vgl.][S.~5]{musterinstitut2023}
Eine Online-Quelle ohne Seitenzahl wird mit Absatz oder Randnummer
zitiert.\autocite[Vgl.][o.\,S.]{musterbehoerde2025}

% --- 3.2 ---
\subsection{Zweiter Analyseschritt}
\label{subsec:analyse2}

Der letzte Absatz eines Unterkapitels zieht eine Folgerung und leitet zum
nächsten über, ohne dessen Ergebnis vorwegzunehmen.
